\documentclass[11pt]{article}

\usepackage[margin=25mm]{geometry}

\usepackage{amsmath,amssymb,fontspec,unicode-math}

\IfFontExistsTF{DejaVu Serif}{\setmainfont{DejaVu Serif}}{\setmainfont{Latin Modern Roman}}

\IfFontExistsTF{Latin Modern Math}{\setmathfont{Latin Modern Math}}{\setmathfont{DejaVu Math TeX Gyre}}

\usepackage{xurl,hyperref}

\setlength{\parindent}{0pt}\setlength{\parskip}{6pt}

\begin{document}

\title{LLMs Encode Harmfulness and Refusal Separately}

\author{ELVIA preliminary AI analysis}\date{21 October 2026}\maketitle

Preliminary assessment: Minor.\par

Independent scientific review remains pending. Author adoption has not been requested.

The inversion evaluator changes the refusal-rate denominator and the extraction path defaults to a mismatched prompt template. The different intervention masks are disclosed and remain a scope qualification.

\section*{Review scope}

Latest-version, claim-focused review of the cited author-paper passages and applicable current author-release paths. The stated counterexamples and scope controls were checked analytically; no author program, old verifier, training run or benchmark was executed.

\section*{Problem 1: Invalid answers change the refusal-rate denominator}

Source: Current arXiv v5 Section 2, PDF p.4; Section 3.5, PDF pp.8–9; src/eval.py, easy\_eval and process\_layer\_evaluation

Assessment: supported; impact: minor

\subsection*{Source claim (paraphrase)}

Section 2 describes refusal over all test examples. `src/eval.py`, `easy\_eval` and `process\_layer\_evaluation`, instead discard invalid inversion labels before computing the saved rate.

\subsection*{Analysis}

For replies No, Certainly and two abstentions, one negative answer among all four is \(1/4\), but the conditional valid-only rate is \(1/2\). The substring branches also accept “unknown” as containing “no” while rejecting uppercase “NO”. These are denominator and parser differences, not estimates of the actual frequency in Figure 4.

\subsection*{Finding and scope}

The released rate is conditional on accepted labels. The original prediction counts are needed to determine the numerical impact.

\subsection*{Proposed correction}

Retain negative, affirmative and invalid counts. Report unconditional rate \(N_{\rm No}/N\), validity, and any separately named valid-only rate. Use an explicit case-insensitive label contract rather than unrestricted substrings.

An explicit local correction is described; author adoption and a corresponding published erratum are not established.

\section*{Problem 2: Direction extraction defaults to the wrong model template}

Source: Current arXiv v5 Section 2, PDF p.3, model-template convention; src/extract\_hidden.py, generate\_directions; src/utils.py, formatInp\_llama\_persuasion

Assessment: supported; impact: minor

\subsection*{Source claim (paraphrase)}

The advertised Qwen/Llama3 extraction path calls the prompt formatter without its model argument; `src/utils.py` defaults that argument to Llama2. Intervention passes the model explicitly.

\subsection*{Analysis}

Thus the same selected model can receive Llama2 markers during extraction and model-specific markers during intervention. This changes the formatted sequence and semantic token positions. The release does not establish the extraction provenance of every cached direction.

\subsection*{Finding and scope}

The public reproduction path has a template mismatch; cached directions and cross-model results remain unverified.

\subsection*{Proposed correction}

Pass the selected model explicitly, use the same template as intervention, and identify instruction positions on the complete tokenized sequence rather than an unverified isolated-suffix length.

An explicit local correction is described; author adoption and a corresponding published erratum are not established.

\section*{Problem 3: Direction and intervention support vary together}

Source: Current arXiv v5 Section 3.5, PDF pp.8–9; Appendix E.1, PDF p.20, Figure 10

Assessment: not\_issue; impact: no\_issue\_found

\subsection*{Source claim (paraphrase)}

Section 3.5 compares harmfulness and refusal interventions; Appendix E.1 explicitly discloses their different token masks.

\subsection*{Analysis}

Writing outcomes as \(Y(v,m)\), comparing \(Y(v_h,m_{\rm inst})\) with \(Y(v_r,m_{\rm all})\) changes both direction and support. It identifies different intervention packages, rather than an isolated direction effect. Appendix E.1 explicitly explains this choice, and Figure 10 retains the same-support refusal controls.

\subsection*{Finding and scope}

This disclosed design choice does not establish a separate error or invalidate all evidence for distinct directions.

\section*{Limits}

Preliminary AI assessment; independent human scientific adjudication remains pending.\par

An analytic counterexample or a release discrepancy does not by itself remeasure the paper’s reported experiments.\par

\section*{Source identity}

arXiv:2507.11878v5 (current author version; updated 2026-07-06T01:46:44Z)

\url{https://arxiv.org/pdf/2507.11878v5}

Reviewed PDF SHA-256: \nolinkurl{e8e3b1972f521631d3a74e7101b293ecaae73ad2059a9f6350f1ef4c3f93577b}

The selected author/publisher source is the same work as the original case, matched by exact title, author identity and official version channel. The exact selected PDF SHA-256 was recomputed in the latest-version pass for Issue \#27. The PDF is a retained author/publisher copy; its original retrieval time is stated only when supported by matching acquisition metadata. File modification time was not substituted.

\end{document}
